\chapter{\abstractname}

Halo exchanges periodically introduce communication overhead in ocean models that use domain decomposition. Nonblocking Message Passing Interface (MPI) operations allow independent computation to be performed between communication initiation and completion, but the degree of communication--computation overlap that can actually be achieved depends on how outstanding communication is progressed.

This thesis investigates how MPI implementation, transport, message size, and progress configuration affect communication--computation overlap in a controlled halo-exchange benchmark.

The Comb benchmark is extended with an adjustable-duration synthetic compute kernel inserted between the initiation and completion of nonblocking communication. Experiments are conducted on the LRZ CoolMUC-4 system, primarily using two nodes with one MPI rank per node and covering 19 message payload sizes. The evaluated configurations include Open MPI over TCP, MPICH over TCP, and Intel MPI over both verbs and TCP. Measurements with and without the injected computation are compared to quantify changes in receive completion time. In addition, a mock backend in Comb that performs no actual network communication is used to estimate processing overhead introduced by the benchmark framework, local memory operations, and timing. Repeated independent runs are used to characterize measurement variability.

The results show that simply using nonblocking MPI communication interfaces does not guarantee effective communication--computation overlap. Without asynchronous progress, Intel MPI exhibits very low measured overlap with both verbs and TCP. Enabling asynchronous progress substantially improves the measured overlap across much of the medium- and large-message range. With progress enabled, Open MPI over TCP shows a sustained decline in measured overlap at larger payload sizes, whereas Intel MPI with both verbs and TCP, as well as MPICH over TCP, maintain relatively high measured overlap over the tested large-message range.

These findings demonstrate that communication--computation overlap must be evaluated for the specific MPI configuration and communication workload. Reductions in receive completion time should be considered together with total runtime and run-to-run variability. The benchmark provides experimental evidence relevant to halo exchanges in ocean models, but does not establish application-level speedups for a complete ocean simulation.
